\documentclass{article}
\usepackage{amsmath}
\usepackage{geometry}
\usepackage{booktabs}
\geometry{margin=1in}

\title{Practical Feed and Speed Selection (Shop-Floor Method)}
\author{Milling Supervisor Notes}
\date{\today}

\begin{document}

\maketitle

\section*{Original Context}
You are selecting feeds and speeds for milling operations based on material, tool, and operation, with practical adjustments from shop-floor experience.

\section*{Step 1: Select Surface Speed (SFM)}

\begin{center}
\begin{tabular}{ll}
\toprule
Material & Carbide SFM Range \\
\midrule
Aluminum (6061) & 800--1200 \\
Aluminum (cast) & 600--900 \\
Mild Steel (1018) & 300--500 \\
Alloy Steel (4140) & 200--400 \\
Stainless Steel (304) & 150--300 \\
Tool Steel (D2) & 100--250 \\
Cast Iron & 400--700 \\
Brass & 500--1000 \\
Titanium (Ti-6Al-4V) & 80--200 \\
Nickel Alloys (Inconel) & 50--150 \\
\bottomrule
\end{tabular}
\end{center}

Adjustments:
\begin{itemize}
\item Poor rigidity: reduce 20--40\%
\item Good coolant: increase slightly
\item Worn machine: reduce
\end{itemize}

\section*{Step 2: Convert to RPM}
\[
RPM = \frac{SFM \times 3.82}{D}
\]
SI equivalent:
\[
N = \frac{1000 \cdot V}{\pi D}
\]
where $V$ is cutting speed in m/min and $D$ in mm.

\section*{Step 3: Chip Load Selection}

Definition:
\begin{itemize}
\item Small tools: diameter $< 12$ mm (0.5 in)
\item Large tools: diameter $\geq 25$ mm (1 in)
\end{itemize}

\begin{center}
\begin{tabular}{lll}
\toprule
Tool Type & Chip Load (in/tooth) & Chip Load (mm/tooth) \\
\midrule
Small endmill & 0.001--0.003 & 0.025--0.075 \\
Medium endmill & 0.002--0.006 & 0.05--0.15 \\
Large face mill & 0.004--0.012 & 0.10--0.30 \\
\bottomrule
\end{tabular}
\end{center}

Shop rule:
\begin{quote}
If chips are dust $\rightarrow$ feed too low. If machine screams $\rightarrow$ feed too high.
\end{quote}

\section*{Step 4: Feed Rate}
\[
Feed = RPM \times Z \times f_z
\]
SI equivalent:
\[
F = N \cdot Z \cdot f_z
\]
where $F$ is mm/min and $f_z$ in mm/tooth.

\section*{Step 5: Practical Adjustments}

\subsection*{Too Aggressive}
\begin{itemize}
\item Vibration or chatter
\item Insert damage
\item Loud cutting noise
\end{itemize}
Action: Reduce feed first, then depth.

\subsection*{Too Conservative}
\begin{itemize}
\item Powder-like chips
\item Rapid tool wear
\item Work hardening
\end{itemize}
Action: Increase feed.

\section*{Special Face Milling Notes}
\begin{itemize}
\item Full-width cuts: reduce feed 20--30\%
\item Interrupted cuts: reduce speed 10--20\%
\item Long tool overhang: reduce both speed and feed
\end{itemize}

\section*{Key Insight}
Feed controls chip formation and heat removal, while speed mainly affects temperature.

\section*{Supervisor Rule}
\begin{quote}
Tables get you close. Chips tell you the truth.
\end{quote}



\section*{Peck Drilling}
Peck drilling is used to break chips, reduce heat, and prevent tool breakage in deep holes.

\subsection*{Basic Rule}
Peck depth ($Q$) depends on diameter $D$:
\begin{itemize}
\item Shallow holes ($< 2D$): no peck needed
\item Moderate depth ($2D-5D$): $Q = 0.5D$ to $1D$
\item Deep holes ($>5D$): $Q = 0.25D$ to $0.5D$
\end{itemize}

Typical formula:
\[
N = \frac{1000 V}{\pi D}, \quad F = N f
\]
where $f$ is feed per rev in mm/rev (SI).

\subsection*{Example}
Drill: 10 mm HSS drill in mild steel
\begin{itemize}
\item Cutting speed: $V = 25$ m/min
\item RPM: $N = \frac{1000\cdot 25}{\pi\cdot 10} \approx 795$
\item Feed per rev: $f = 0.15$ mm/rev
\item Feed rate: $F \approx 120$ mm/min
\item Peck depth: $Q = 5$ mm (0.5D)
\end{itemize}

Guidelines:
\begin{itemize}
\item Reduce peck as hole deepens
\item Use coolant or air blast for chip evacuation
\item Retract fully every 3--5 pecks in deep holes
\end{itemize}

\section*{Entry Methods for Pocket Machining}

Entry strategy prevents tool damage and improves tool life.

\subsection*{1. Plunge Entry}
\begin{itemize}
\item Direct vertical entry
\item Use only for center-cutting tools
\item Feed: 30--50\% of normal cutting feed
\end{itemize}

Example:
12 mm endmill, normal feed = 800 mm/min
\begin{itemize}
\item Plunge feed: 300--400 mm/min
\end{itemize}

\subsection*{2. Ramp Entry}
\begin{itemize}
\item Tool enters at angle (1--3 degrees typical)
\item Safer than plunge
\item Use for most pocketing operations
\end{itemize}

Ramp angle relation:
\[
\tan(\theta) = \frac{\text{depth}}{\text{horizontal travel}}
\]

Example:
Depth = 10 mm, angle = 2$^\circ$
\begin{itemize}
\item Horizontal distance $\approx 286$ mm
\end{itemize}

\subsection*{3. Helical Entry}
\begin{itemize}
\item Spirals into material
\item Best load distribution
\item Preferred for deep pockets
\end{itemize}

Typical parameters:
\begin{itemize}
\item Helix diameter: 1.2--1.5$\times$ tool diameter
\item Pitch: 0.5--1.0 mm per revolution
\end{itemize}

Example:
16 mm endmill
\begin{itemize}
\item Helix dia: 20--24 mm
\item Pitch: 0.8 mm/rev
\end{itemize}

\subsection*{4. Pre-drill Entry}
\begin{itemize}
\item Drill hole first, then mill
\item Useful for hard materials or large depths
\end{itemize}

\section*{Supervisor Insight}
\begin{quote}
Never force a tool straight into material unless it is designed for it. Entry strategy determines tool life more than cutting speed.
\end{quote}

\end{document}
